% SPDX-License-Identifier: CC-BY-4.0
\documentclass[11pt,letterpaper]{article}

\usepackage[T1]{fontenc}
\usepackage{amsmath,amssymb}
\usepackage{tikz}
\usepackage{pgfplots}
\pgfplotsset{compat=1.18}
\usepackage{algorithm}
\usepackage{algpseudocode}
\usepackage{booktabs}
\usepackage{enumitem}
\usepackage[hidelinks]{hyperref}

\newcommand{\N}{\mathbb{N}}
\newcommand{\Z}{\mathbb{Z}}
\newcommand{\Q}{\mathbb{Q}}
\newcommand{\R}{\mathbb{R}}
\newcommand{\C}{\mathbb{C}}
\newcommand{\Map}[3]{#1 \colon #2 \to #3}
\newcommand{\Mod}[1]{#1\text{-}\mathrm{Mod}}
\newcommand{\GroupHom}{\mathrm{Hom}_{\mathbf{Grp}}}
\newcommand{\RingHom}{\mathrm{Hom}_{\mathbf{Ring}}}

\newtheorem{theorem}{Theorem}
\newtheorem{lemma}[theorem]{Lemma}
\newtheorem{definition}[theorem]{Definition}
\newtheorem{remark}[theorem]{Remark}
\newenvironment{proof}{\noindent\textit{Proof.}}{\hfill\(\square\)}

\title{Generalized Collatz Library}
\author{ascendnoosphere}
\date{\today}

\begin{document}
\maketitle

\begin{abstract}
A positive common-base divisor and a sufficiently divisible offset force every
nonzero signed raw orbit to escape in magnitude and ultimately have the
offset's sign. A complete phase sweep gives an exact arithmetic progression at
fixed raw times. Whole-base order permits composite bases and nonunit
residuals. Sign-qualified intermediate bounds recover the full raw sequence
with an explicit finite budget. Zero is the only periodic point. A primitive
cycle at the preceding offset threshold shows that the sufficient threshold
cannot uniformly be lowered by one.
\end{abstract}

\section{Parameters, clocks and scope}

% These are the exact common hypotheses of the formal escape statements.
\begin{definition}
Let \(t,u,v,a\in\N\) and \(s\in\Z\) satisfy
\[
t\ge2,\qquad 0<u<v,\qquad \gcd(u,v)=1,\qquad a\ge u+v,\qquad s\ne0.
\]
Put \(Q=t^u\), \(D=t^v>0\), \(C=t^a s\), and define the raw map
\[
F(n)=\begin{cases}n/D,&D\mid n,\\Qn+C,&D\nmid n.\end{cases}
\]
Division has priority, including at zero. Both seeds and offsets may be
negative; the divisor in this theorem is positive.
Write \(G_v(t)=\sum_{j=0}^{v-1}t^j\).
\end{definition}

The positive-offset common-base regimes are organized in the source guide
\cite{QcdPowerPhase}. Here we prove the super-offset branch \(a\ge u+v\),
with a signed extension. Neither primality of \(t\), coprimality of \(s,t\),
nor \(t\nmid s\) is imposed. The sub-offset, intermediate and critical regimes
are separate obligations. The broader signed-divisor family is not identified
with this branch.

For later use define a normalized event and its actual raw clock by
\[
\tau(x)=
\begin{cases}2,&D\mid F(x),\\1,&D\nmid F(x),\end{cases}
\qquad E(x)=F^{\tau(x)}(x).
\]
On the normalized states considered below, this performs one affine instruction
followed by every immediately available division, whose number is zero or one.
It is not a first-division map or an odd shortcut.

\section{Whole-base order and initial divisions}

% Whole-base order uses integer cancellation, not prime valuations.
\begin{definition}
A nonzero integer has whole-base order \(k\) at base \(t\) if
\[
n=t^k w,\qquad t\nmid w.
\]
The residual \(w\) may share primes with \(t\). For example, two has order zero
at base four, and \(-32=4^2(-2)\) has order two.
\end{definition}

% Strict magnitude descent constructs the witness and terminates the initial prefix.
\begin{lemma}[Initial divisions]\label{Normalize}
A nonzero seed \(n\in\Z\) has a legal raw division prefix of length
\(r_0\le |n|\) reaching \(x=t^k w\), where
\[
0\le k<v,\qquad t\nmid w,\qquad |x|\le|n|.
\]
In particular \(D\nmid x\).
\end{lemma}
\begin{proof}
If \(t\mid n\), division by \(t\) preserves nonzeroness and strictly decreases
magnitude, since \(t\ge2\). Repeating therefore terminates at a nonzero \(w\)
not divisible by \(t\), and produces \(n=t^m w\). Cancellation shows the
order is exact: an additional factor \(t\) in \(n\) would divide \(w\).
Since \(|w|\ge1\), one has \(t^m\le|n|\), and \(m\le t^m\) gives
\(m\le|n|\).

Write \(m=vr_0+k\), \(0\le k<v\). Then
\(n=D^{r_0}t^k w\); the first \(r_0\) raw instructions are legal divisions and
reach \(x=t^k w\). Their quotients are exact, so \(|x|\le|n|\), and
\(r_0\le m\le|n|\). If \(D\mid x\), then \(t^{k+1}\mid x\), because
\(k+1\le v\). Canceling \(t^k\) would give \(t\mid w\), a contradiction.
\end{proof}

\section{The deterministic phase sweep}

% The high offset removes cancellation in the leading whole-base factor.
\begin{lemma}[One legal normalized event]\label{Event}
If \(x=t^k w\), \(k<v\), and \(t\nmid w\), then
\[
Qx+C=t^{k+u}\bigl(w+s\,t^{a-u-k}\bigr)
\]
has exact whole-base order \(k+u\). The normalized event uses zero divisions
when \(k+u<v\), and one division otherwise. Its next phase is
\((k+u)\bmod v\), with residual \(w+s\,t^{a-u-k}\).
\end{lemma}
\begin{proof}
The exponent \(a-u-k\) is at least one, since \(a\ge u+v\) and \(k<v\).
Thus the correction is divisible by \(t\), while \(w\) is not; their sum
is not divisible by \(t\) and is nonzero. This proves exact order \(k+u\),
without any assumption on the sign or prime factors of \(w,s\).
Because \(0\le k+u<2v\), division by \(t^v\) is either unavailable or occurs
exactly once. In the latter case the quotient has order \(k+u-v<v\).
The initial normalized state is not divisible by \(D\), so the first raw
instruction really is affine. Both subsequent possibilities are legal and
exhaustive.
\end{proof}

% Cancellation modulo v closes all phases rather than selecting favorable ones.
\begin{lemma}[Phase permutation and raw clock]\label{Permutation}
Let \(k_0=k\), \(k_i=(k+iu)\bmod v\). For \(0\le i<v\), these phases
permute \(\{0,\ldots,v-1\}\). In the first \(i\) normalized events the
division count and raw clock are
\[
e_i=\left\lfloor\frac{k+iu}{v}\right\rfloor,\qquad T_i=i+e_i.
\]
In particular a sweep of \(v\) events has \(u\) divisions and raw length
\(P=u+v\).
\end{lemma}
\begin{proof}
If \(k_i=k_j\), then \(v\mid u(i-j)\). Coprimality gives
\(v\mid(i-j)\); since \(0\le i,j<v\), this forces \(i=j\).
An injective map of the finite \(v\)-element phase set to itself is a
permutation. This argument requires \(\gcd(u,v)=1\), not that the base be prime.

At each step Lemma~\ref{Event} increases the phase by \(u\), subtracting
\(v\) exactly when a division occurs. Telescoping gives
\(k_i=k+iu-ve_i\), with \(0\le k_i<v\), so \(e_i\) is the stated quotient.
Each event has one affine instruction and its zero or one division, hence
\(T_i=i+e_i\). At \(i=v\), \(k_v=k\) and
\(e_v=\lfloor(k+vu)/v\rfloor=u\), proving the final claim.
\end{proof}

% The finite geometric kernel supplies a nonzero, phase-preserving translation.
\begin{theorem}[Exact same-seed raw ray]\label{Ray}
For the normalized state \(x=t^k w\), put
\[
\beta_k=s\,t^{a-u-v+1+k}G_v(t),\qquad P=u+v.
\]
Then \(F^{jP}(x)=x+j\beta_k\) for every \(j\in\N\).
For the initial prefix of Lemma~\ref{Normalize},
\[
F^{r_0+jP}(n)=x+j\beta_k
\]
at the same original signed seed.
\end{theorem}
\begin{proof}
Each normalized event adds \(s\,t^{a-u-k_i}\) to the residual.
After \(v\) events the phase is again \(k\), and the residual addition is
\[
s\sum_{i=0}^{v-1}t^{a-u-k_i}
=s\,t^{a-u-v+1}\sum_{i=0}^{v-1}t^{v-1-k_i}
=s\,t^{a-u-v+1}G_v(t).
\]
All exponents are nonnegative, and Lemma~\ref{Permutation} justifies the
last equality over the complete phase set. Multiplication by \(t^k\)
gives \(F^P(x)=x+\beta_k\).

The exponent \(a-u-v+1+k\) is at least \(k+1\). Therefore
\(\beta_k/t^k\) is an integer divisible by \(t\), and
\(w+j\beta_k/t^k\) remains not divisible by \(t\) for every \(j\).
All translated states have the same phase \(k\) and are nonzero.
The same legal sweep applies at each endpoint, so induction proves the
ray formula. Composing the legal initial prefix proves its same-seed version.
\end{proof}

\section{Full raw recovery and the signed escape budget}

% The sign hypothesis prevents cancellation between the affine state and the offset.
\begin{lemma}[Aligned intermediate bound]\label{Aligned}
If a normalized state \(x=t^k w\) satisfies \(sx\ge0\), then for every
raw future state \(z=F^r(x)\),
\[
t^k|z|\ge|x|.
\]
If \(sx>0\), every raw future state also has strict sign \(sz>0\).
\end{lemma}
\begin{proof}
After \(i\) events, the residual has the form
\[
w+sB_i,\qquad
B_i=\sum_{j<i}t^{a-u-k_j}\ge0.
\]
The phase is \(k_i\ge0\). Since \(sw\ge0\), adding \(sB_i\) cannot reduce
the magnitude of \(w\), so every event endpoint has magnitude at least
\(|w|\). When an event has an interior raw state, it is the affine value
before its division. The multiplier \(Q=t^u\ge2\) is positive and the
offset has sign \(s\); under alignment its magnitude is at least that
of the event input, and hence at least \(|w|\).

Every raw time lies in an event block because its clock is positive.
For completeness, accumulated clocks increase by one or two; induction on
the raw time either increases its interior offset or passes to the next block.
This is the positive-clock coverage principle of \cite{AscendCompressionII}.
Consequently \(|z|\ge|w|\) at every raw time and
\(t^k|z|\ge t^k|w|=|x|\). If the initial product with \(s\) is positive,
affine instructions retain that sign by addition of a same-sign offset,
and division by positive \(D\) retains it as well.
\end{proof}

The hypothesis cannot be discarded. For \((Q,C,D)=(2,8,4)\), a raw sweep
from \(-5\) is
\[
-5\longrightarrow-2\longrightarrow4\longrightarrow1.
\]
Its phase is zero, and the interior magnitude two is smaller than five.
Even the coarser factor \(D^u=4\) fails at the last magnitude one.
The ray nonetheless has translation six and subsequently aligns.

% Integer ray crossing precedes application of the sign-qualified lower bound.
\begin{theorem}[Effective full-sequence escape]\label{Escape}
For every nonzero \(n\in\Z\) and \(M,r\in\N\), define
\[
B(n,M)=|n|+\bigl(|n|+1+MD^u\bigr)(u+v).
\]
If \(B(n,M)\le r\), then \(|F^r(n)|>M\).
For every \(r\ge B(n,0)\), one also has \(sF^r(n)>0\).
\end{theorem}
\begin{proof}
Normalize \(n\) as in Lemma~\ref{Normalize}. Put \(K=t^k\).
One has \(1\le K\le D^u\), because \(k<v\) and \(u\ge1\).
The integer \(\beta_k\) is nonzero, has the sign of \(s\), and
\(|\beta_k|\ge1\). Choose
\[
J=|x|+1+MK,\qquad y=x+J\beta_k.
\]
If \(\beta_k>0\), then \(y\ge-|x|+J>MK\).
If \(\beta_k<0\), then \(y\le|x|-J<-MK\).
Thus \(sy>0\) and \(|y|>MK\). By Theorem~\ref{Ray}, this state occurs
at raw time \(r_0+JP\), and retains phase \(k\).

Apply Lemma~\ref{Aligned} from \(y\), now that its sign hypothesis holds.
Every later raw state \(z\) satisfies \(K|z|\ge|y|>MK\);
since \(K>0\), it follows that \(|z|>M\).
The time used obeys
\[
r_0+JP\le |n|+\bigl(|n|+1+MD^u\bigr)P=B(n,M),
\]
using \(r_0\le|n|\) and \(|x|\le|n|\). This proves the assertion for
every raw time beyond the displayed budget, not only sweep endpoints.
For sign alone take \(M=0\); the same crossing and the strict-sign part of
Lemma~\ref{Aligned} retain \(sz>0\) at all later raw times.
\end{proof}

% Escape excludes every positive return while zero retains its legal division loop.
\begin{theorem}[Periodic points]\label{Periodic}
Zero is the only periodic point of \(F\), and its primitive period is one.
\end{theorem}
\begin{proof}
Zero takes division and is fixed. Suppose a nonzero \(n\) had a positive
return time \(h\). Determinism propagates this return to every multiple
\(jh\). Choose a multiple beyond \(B(n,|n|)\). Theorem~\ref{Escape} then gives
\(|F^{jh}(n)|>|n|\), contradicting \(F^{jh}(n)=n\).
The least positive return at zero is one.
\end{proof}

\section{Threshold and concrete arithmetic}

% A primitive legal boundary cycle refutes a uniform lowering of the threshold.
\begin{theorem}[Boundary obstruction]\label{Boundary}
At \(a=u+v-1\), the parameter choice
\[
(t,u,v,a,s)=(2,1,2,2,1)
\]
has the primitive positive raw cycle
\[
1\longrightarrow6\longrightarrow16\longrightarrow4\longrightarrow1.
\]
Hence the uniform sufficient threshold cannot be replaced by \(a\ge u+v-1\).
\end{theorem}
\begin{proof}
Here \(Q=2,C=4,D=4\). One and six are not divisible by four, so their
successors are \(2\cdot1+4=6\) and \(2\cdot6+4=16\).
Sixteen and four take legal divisions to four and one.
None of the first three positive times returns to one, proving primitiveness.
At this threshold the leading residual correction in Lemma~\ref{Event}
need not contain a factor \(t\). Indeed the second affine value has extra
divisibility and permits two successive divisions; the higher-offset phase
argument no longer applies.
\end{proof}

For \((t,u,v,a,s)=(3,2,3,5,1)\), one has \(Q=9,C=243,D=27\) and
\[
1\longrightarrow252\longrightarrow2511\longrightarrow93
\longrightarrow1080\longrightarrow40.
\]
There are three affine instructions and two divisions, so \(P=5\).
The kernel is \(\beta_0=3(1+3+9)=39\); the next endpoint is \(79\).

For \((t,u,v,a,s)=(4,1,3,4,2)\), \(Q=4,C=512,D=64\), and
\[
32\longrightarrow640\longrightarrow10\longrightarrow552\longrightarrow2720.
\]
Here \(32=4^2\cdot2\), so \(k=2\) with a nonunit residual.
The offset residual two is also nonunit modulo four.
The raw clock is four and \(\beta_2=2\cdot4^3(1+4+16)=2688\).
This explicitly tests why whole-base order must not be replaced by a prime
valuation or a unit-residual assumption.

\section{Reduction certificate and finite computation}

The original obligation is the orbit classification of every signed seed in
the stated super-offset parameter family. The replacement data are a finite
legal division prefix, a phase \(k<v\), an explicitly evaluated finite kernel,
and a nonzero signed arithmetic progression. The initial prefix terminates
by a strictly decreasing integer magnitude. The high-offset condition removes
leading-order cancellation, and the finite coprime sweep closes every phase.
The sign-qualified intermediate bound and positive-clock coverage reconstruct
every raw time with the displayed budget. Thus this is a proved solvable
parameter subclass; no infinite-orbit oracle or unproved uniform estimate
has been hidden in the replacement data. The budget is effective in integer
values, not claimed polynomial in their bit lengths.

The Go implementation replays two raw sweeps, checking every branch, phase,
division count, clock and translation endpoint. The main grid uses
\[
\begin{aligned}
2&\le t\le9,&2&\le v\le8,&0&<u<v,\quad\gcd(u,v)=1,\\
u+v&\le a\le u+v+2,&-8&\le s\le8,\quad t\nmid s.
\end{aligned}
\]
It uses ordinary nonzero seeds in \([-128,128]\) and all phases \(t^k w\),
\(k<v\), with \(-8\le w\le8\), \(t\nmid w\). Its reference count is
2,186,784 two-sweep certificates, not distinct orbits.
Separate tests include offsets \(\pm3\cdot2^{204}\), opposite signs,
a divisible offset residual and the exact boundary cycle.
Sign-qualified magnitude checks are deliberately skipped before alignment.

The public finite checker refuses parameters outside its explicit contract:
\(2\le t\le9\), \(0<u<v\le8\), coprimality,
\(u+v\le a\le256\), and \(0<|s|\le8\).
It does not require \(t\nmid s\).
Its reserved seed bit limit is
\(\max(8,(v-1)\operatorname{BitLen}(t)+4)\);
nil, unconfigured, mismatched, oversized or zero seed requests are rejected
before replay. These limits restrict computation, not the mathematical
theorem. Disjoint reusable buffers and cached powers provide zero warmed
allocations on the four tested configurations; no arbitrary-size performance
assertion follows. Finite replay is distinct from the universal proof.

The six formal proof modules are
\begin{itemize}
\item \texttt{CommonBaseHighOffset.lean}: whole-base order and legal events;
\item \texttt{CommonBaseHighOffsetKernel.lean}: complete phase kernel;
\item \texttt{CommonBaseHighOffsetClock.lean}: actual raw clocks and rays;
\item \texttt{CommonBaseHighOffsetNormalization.lean}: the legal initial prefix;
\item \texttt{CommonBaseHighOffsetEscape.lean}: full raw magnitude and sign;
\item \texttt{CommonBaseHighOffsetOutcomes.lean}: periodicity and threshold.
\end{itemize}

\bibliographystyle{amsalpha}
\bibliography{CommonBaseHighOffset}
\end{document}
